% =========================================================================
% CHAPTER 10 — DISCUSSION
% =========================================================================

\chapter{Discussion}
\label{ch:discussion}

\noindent
{\rmfamily\itshape
\begin{tabular}[t]{@{}l@{\hspace{0.3cm}}p{12.5cm}@{}}
\textls[50]{\textbf{Addressed:}} &
\textls[50]{Interpretation · Expectations Revisited · Mechanism \& Boundary Conditions ·
Theoretical, Practical \& Sustainability Implications · What the Study Can and Cannot Claim}
\end{tabular}
}

\vspace{0.5cm}

Chapter~\ref{ch:results} reported what the coded responses show. This chapter asks what they mean. It returns to the expectations that framed the user study, reads the results against them, and states plainly where the evidence supports the thesis, where it complicates it, and where it cannot speak. The reading is careful for a reason. With twenty-five participants, the study cannot prove a general claim. What it can do is show whether the patterns it was built to look for actually appear, and the most useful result of a first study is often not a confirmation but a sharpened understanding of where the approach helps and why.

% -------------------------------------------------------------------------
\section{Reading the Expectations Against the Study as Built}
\label{sec:disc-translate}

The expectations for the evaluation, E3.1 and E3.2, were written early, for a design that later changed: three conditions of increasing experiential richness and a think-aloud protocol became two conditions and a written self-administered survey (Chapter~\ref{ch:userstudy}). They therefore need to be read in translated form. E3.1 --- that greater experiential richness would show in \emph{how} participants articulate their understanding rather than in a comprehension score --- is read as a comparison between the two conditions. E3.2 --- that the perceptual encodings would produce interpretations the protocol could surface, matching the intended reading or not --- is read as applying to what the written responses reveal.

This translation is a limitation, and it is stated as one. But it does not weaken the expectations, because E3.1 was cautious in exactly the way that turned out to matter. It did not predict that the experiential condition would raise comprehension. It predicted a difference in articulation, not in score. That caution turns out to matter, because it is precisely what the results show.

% -------------------------------------------------------------------------
\section{E3.1: The Cautious Expectation Holds, the Naive One Would Not}
\label{sec:disc-e31}

The clearest single result is that the experiential condition did not produce a uniform gain in comprehension. Across the three scenes the direction of the difference was not the same (Figure~\ref{fig:grasped}): level in Scene~I, ahead for the experimental condition in Scene~II, ahead for the control in Scene~III. A study that had expected experiential presentation to lift understanding everywhere would be embarrassed by this pattern. The present study is not, because it never expected that.

This is the heart of the matter. Had the expectation been the naive one --- that a richer, moving representation simply produces better comprehension --- the data would have refused it. Comprehension scores did not rise with experiential richness. What changed instead was the character of the understanding, and where it was reachable at all. This is exactly the shape of E3.1: a difference in how the system is grasped and articulated, not a general gain in a score. The warning that animated visualisations can look effective only because they smuggle in information the static version lacks \parencite{tversky2002} is the reason the study held its two conditions informationally equal; the result is that no such artefact inflates the experiential condition, and the flat comparison in Scene~I is honest rather than disappointing.

% -------------------------------------------------------------------------
\section{Where the Difference Lives, and Why}
\label{sec:disc-mechanism}

If the effect is not a uniform lift, the interesting question is where it appears and why. The three scenes answer it together.

Scene~II concerns recurrence: whether the system passes back through states it has already occupied. A still scatter cannot show this, because without an ordering in time there is nothing to reveal a return. The results match exactly. No control participant was coded as grasping recurrence, against roughly a quarter of the experimental group, with partial understanding also rising (Chapter~\ref{ch:results}). The mechanism is not inferred; the participants state it themselves. Every response marked \textsc{Interaction-aided} --- attributing understanding to pausing, rewinding or isolating part of the display --- came from the experimental condition, and every response marked \textsc{Static-limit-named} --- naming the absence of a temporal ordering as the reason recurrence could not be judged --- came from the control condition (Figure~\ref{fig:emergent}). One group says the motion let them see the return; the other says the still image did not. These are the two faces of the same claim, spoken from opposite sides of the design.

A third emergent code cuts against the experiential condition, and is reported here for the same reason the Scene~III reversal is. \textsc{Anomaly-confusion}, marking difficulty with the usual/unusual encoding, appeared five times in the experimental condition against once in the control (\S\ref{sec:res-emergent}). Principle~P3 assigns anomaly to texture and agitation, on the grounds that irregularity draws attention without a legend \parencite{bertin1983, ware2021}. These responses suggest that drawing attention is not the same as conveying meaning: participants noticed the disturbance and could not say what it signified. Pre-attentive salience appears to be sufficient for detection and insufficient for interpretation --- a limit of P3 as stated, and a specific target for revision.

Scene~III is the test that makes this credible, because it runs the other way. Here the control condition did better, reaching every participant against roughly three-quarters of the experimental group. The reason is instructive rather than awkward. Scene~III asks whether the agreement between three models changes over time, and a static figure shows the whole time series at once, laying the convergence out to be read in a single view. The experiential version asks the participant to build that same reading by moving through the controls, and a few did not complete it. A second explanation cannot be separated from the first at this design: Scene~III came last for every participant, and the experiential version asks for sustained attention precisely where attention is least likely to remain. The fixed scene order (\S\ref{sec:eval-limitations}) means the two accounts are confounded, and the study cannot adjudicate between them. Under either reading, the lesson is not that the experiential form failed. It is that experiential presentation is not free. It helps most where the structure is invisible to a still figure, as in Scene~II, and can cost attention where a static figure already presents the structure plainly, as in Scene~III. An approach that admits this is stronger than one that claims a uniform advantage, because it can say \emph{where} it applies.

% -------------------------------------------------------------------------
\section{E3.2 and the Streamfunction Misconception}
\label{sec:disc-e32}

The most substantial finding of the study is a misreading, and E3.2 anticipated exactly this kind of thing: that non-specialists would interpret the representation in ways worth surfacing, whether or not the interpretation matched the intended one. Fourteen of twenty-five participants read the height of the streamfunction curve in Scene~I as the speed of the water, rather than as the cumulative transport it represents (Chapter~\ref{ch:results}). The curve invites the reading. Its height is the most salient thing on the display, and height reads naturally as amount or speed; the correct reading, that direction and strength live in the curve's slope, requires an inference that a non-specialist does not spontaneously make.

This misreading is not domain-specific, and the literature of Chapter~\ref{ch:literature} already names it. The streamfunction is a cumulative integral of transport and its slope is the local rate (Equation~\ref{eq:velocity}); reading its height as speed is precisely the stock--flow confusion that \textcite{sterman2008} and \textcite{cronin2009} document in mathematically competent adults asked to relate a rate to the level it accumulates. What the present study adds is that the confusion survives translation into a different domain and a different visual idiom: participants met an oceanographic curve rather than a stock-and-flow diagram, and made the same substitution. The finding therefore does more than record a local design flaw. It supplies a case in which a known limit of cognitive architecture reappears in the exact place the cognitive-science strand of this thesis predicted it would --- which turns the Scene~I result from an anomaly into a confirmation.

Two further patterns turn this from an error into a finding. The first is that the misconception was more frequent in the experimental condition, not less (Figure~\ref{fig:misconception}). At first sight this contradicts the thesis. It does not, once the second pattern is placed beside it: the misconception fell sharply with data-visualisation familiarity, from roughly three-quarters of the less familiar participants to a third of the more familiar. The experiential display shows moving particles, and moving particles make the velocity reading more immediate, not less. The perceptual richness that helped participants see recurrence in Scene~II here works against them, making the seductive reading easier for exactly those with the fewest tools to correct it. This is not a failure of the approach so much as a lesson about it: perceptual salience is powerful in both directions, and a representation that makes one structure vivid can make an adjacent misreading vivid too.

That the two conditions also differed in mean data-visualisation familiarity means condition and familiarity cannot be fully separated at this sample size (Chapter~\ref{ch:results}), and the joint reading is therefore exploratory. But the direction is clear enough to carry a design consequence, taken up in Chapter~\ref{ch:conclusion}: the cumulative nature of the curve is the thing a future iteration must make perceptible, rather than leaving the viewer to supply it.

The misconception also gives E3.2 its confidently-wrong edge. \textcite{fischer2020} describe cases where high confidence attaches to an incorrect reading, and Scene~I contains them: four participants rated their confidence at seven or above while reading the curve only partially, each of them carrying the streamfunction misconception (Figure~\ref{fig:confidence}). Confidence, in this scene, does not track comprehension. The people most sure of their reading were among those who had misread the instrument, which is precisely why the study kept perceived and demonstrated understanding apart rather than letting a confidence rating stand in for a grasp.

% -------------------------------------------------------------------------
\section{What the Study Cannot Claim}
\label{sec:disc-limits}

Three limits bound these readings, and naming them is part of the argument rather than an apology appended to it.

The sample is small and was recruited by convenience, so nothing here generalises to a population and the exploratory tests carry no inferential weight. The Scene~III comparison, where the control did better, is as real as the Scene~II one, and the study reports both rather than the flattering one alone. The confound between condition and data-visualisation familiarity means the two panels of the misconception result are not independent: the finding is a direction to test, not a settled effect. And the coding rests on a single researcher supported by an automated first pass, so no inter-coder reliability can be reported (Chapter~\ref{ch:userstudy}). None is fatal to a first study; each is a reason the next should be larger, balanced on prior familiarity, independently coded, and counterbalanced on scene order.

% -------------------------------------------------------------------------
\section{Implications}
\label{sec:disc-implications}

\paragraph{Theoretical.} The study supplies a case in which the stock--flow confusion of \textcite{sterman2008} reappears in a scientific visual idiom rather than a textbook diagram (\S\ref{sec:disc-e32}), and a case in which the equivalence principle of \textcite{tversky2002} does real work: with information held equal, the advantage of animation does not appear uniformly, and its absence in Scene~I is informative rather than disappointing. It also qualifies the animated-uncertainty literature \parencite{hullman2015hops, kale2019hops}: Scene~III animated uncertainty and did not outperform the static baseline, which suggests the benefit of animating uncertainty may depend on whether the static alternative already affords a whole-series view.

\paragraph{Practical.} For anyone designing a representation of an invisible system, the result is a selection rule rather than an endorsement: motion earns its cost where the structure is temporal and unavailable to a still figure, and adds attentional load without comprehension where a static figure already lays it out. Perceptual salience must also be audited in both directions --- the same particle field that made counter-flow visible in Scene~I made the velocity misreading more available, and did so most for the participants least equipped to correct it.

\paragraph{Sustainability and communication.} That comprehension does not follow perceived clarity extends \textcite{kause2020} into an experiential format, and bears on how climate material is prepared for non-specialist audiences: an artefact evaluated on preference or engagement alone can be optimised in the wrong direction. Working memory is the binding constraint \parencite{sweller1988}, and a representation that spends it on decoration rather than on structure spends it badly.

% -------------------------------------------------------------------------
\section{What This Means for H-A and H-B}
\label{sec:disc-umbrella}

The umbrella hypotheses framed the whole enquiry: that conventional representations are perceptually suboptimal for non-specialists because they treat complexity as an object to be read rather than a dynamic to be inhabited (H-A), and that machine-learning-derived, experiential representations can support understanding by rendering structure perceptually available (H-B). Twenty-five participants cannot prove either. Both find their clearest support in Scene~II, where the static representation left recurrence simply unreachable and participants said so in their own words, and in those who attributed their understanding to acting on the display.

The more valuable contribution is the boundary the study draws around H-B. Experiential presentation is not uniformly better: it helps where structure is temporal and invisible to a still figure, it can hinder where a static figure already lays structure out, and its perceptual power can deepen a misreading as readily as a reading. The thesis emerges not as proven but as sharpened --- the question is no longer whether experiential representation helps, but for which structures, for which viewers, and at what perceptual cost.